\documentclass[../main.tex]{subfiles}

\graphicspath{{../image/}{../figure/}}

\begin{document}

\chapter{平衡律}\label{chapter:ch01-平衡律}

本章阐述平衡律的一般数学理论，旨在为经典连续介质物理中众多平衡律提供一个统一框架；这些平衡律由诸如质量、动量、能量等所谓“广延量”所满足。平衡律所处的环境空间为 $\mathbb{R}^{k}$，典型点记为 $X$。在连续介质物理的应用中，静力学语境下 $\mathbb{R}^{k}$ 表示一维、二维或三维物理空间；动力学语境下表示二维、三维或四维时空。

一般平衡律将通过其原始表述引入：公设任一域内广延量的生成由通过边界的通量所平衡；随后将其归约为一个场方程。正是这一归约使连续介质物理在数学上可处理。我们将证明，场方程的散度形式在坐标变换下保持，且原始形式的平衡律可从场方程恢复。本章讨论的性质完全由场方程的散度形式导出，因而同等适用于支配平衡与演化的平衡律。

平衡律方程组的场方程将与本构方程联立——后者把通量密度与生成密度联系到一个状态向量——从而得到一个封闭的、散度形式的拟线性一阶偏微分方程组。

我们将证明，可对称化的平衡律方程组配有伴随平衡律，它们被光滑解自动满足，但不必被弱解满足。弱解的容许性问题将被提出。

我们将考虑带激波阵面或弱阵面的解，即状态向量场或其导数在余维 $1$ 流形上经历跳跃间断的解。

为带激波阵面的解提供自然框架的 $BV$ 函数理论将被概述，$BV$ 解的几何结构将被描述。高度振荡弱解将被构造，容许性条件的稳定化作用将被初步指示。

环境空间取 Euclid 空间，故宜采用矩阵记号，代价是掩盖场的张量本性。符号 $\mathbb{M}^{r\times s}$ 此后恒表示 $r\times s$ 矩阵组成的向量空间，$\mathbb{R}^{r}$ 与 $\mathbb{M}^{r\times 1}$ 等同。$\mathbb{R}^{r}$ 中 $(r-1)$ 维 Hausdorff 测度记为 $\mathscr{H}^{r-1}$。其他标准记号包括：$\mathbb{S}^{r-1}$ 表示 $\mathbb{R}^{r}$ 中的单位球面，$\mathscr{B}_{\rho}(X)$ 表示以 $X$ 为心、半径为 $\rho$ 的球；特别地，$\mathscr{B}_{\rho}\triangleq\mathscr{B}_{\rho}(0)$。

\subfile{Chapter-01/section-01.tex}
\subfile{Chapter-01/section-02.tex}
\subfile{Chapter-01/section-03.tex}
\subfile{Chapter-01/section-04.tex}
\subfile{Chapter-01/section-05.tex}
\subfile{Chapter-01/section-06.tex}
\subfile{Chapter-01/section-07.tex}
\subfile{Chapter-01/section-08.tex}
\subfile{Chapter-01/section-09.tex}
\subfile{Chapter-01/section-10.tex}

\end{document}
